\documentclass[12pt, twoside]{article}
\input{../preamble}
\setlength{\parskip}{0.3\baselineskip}

\fancyhead[LE]{\thepage}
\fancyhead[RO]{Markscheme}
\fancyhead[LO]{La Scuola d'Italia / Dr.\ Huson / IB Math 12 \\* 15 September 2026}

\begin{document}
\subsubsection*{1.3 Classwork: Mean, Median and Mode --- Markscheme}

\emph{Data sources: Tax Foundation, \emph{VAT Rates in Europe 2026} and
\emph{State and Local Sales Tax Rates 2026} (state rate only). Both sets are
real 2026 figures.}

\begin{enumerate}[itemsep=0.3cm]

%--- Problem 1: EU VAT, right skew [8 marks] ---
\item[1.] [Maximum mark: 8]

Twelve standard VAT rates, sorted:
$$18,\; 19,\; 19,\; 20,\; 20,\; \underline{21},\; \underline{21},\; 21,\; 22,\; 23,\; 25,\; 27$$
$n = 12$, $\sum x = 256$.

\textbf{(a)} $\bar{x} = \dfrac{256}{12} = \dfrac{64}{3} = 21.333\ldots = 21.3$
(3 s.f.) \hfill \textbf{M1A1}

\emph{M1 for a sum of the twelve values divided by 12, by hand or on the GDC;
A1 for 21.3. Accept $\frac{64}{3}$ or $21\frac{1}{3}$. Do not accept 21 --- that
is the median and the mode.}

\emph{\textbf{GDC:} \texttt{menu} $\rightarrow$ \texttt{Statistics}
$\rightarrow$ \texttt{Stat Calculations} $\rightarrow$
\texttt{One-Variable Statistics}.}

\textbf{(b)} With $n = 12$ the median is the mean of the 6th and 7th values:
$\dfrac{21 + 21}{2} = 21$ \hfill \textbf{M1A1}

\emph{M1 for an ordered list or for identifying the middle pair; A1 for 21.
A candidate who takes a single middle value still reaches 21 here --- award both
marks, but the return should say that with an even $n$ the two middle values
are averaged.}

\textbf{(c)} Mode $= 21$ (Spain, Netherlands and Belgium).
\hfill \textbf{A1}

\emph{Do not accept 3, which is the frequency of the mode.}

\textbf{(d)} The distribution is \textbf{skewed right} (positive skew).
\hfill \textbf{A1}

Because $\bar{x} = 21.3 > \text{median} = 21$, and the dot plot has a longer
tail on the high side while most values bunch between 19 and 22.
\hfill \textbf{R1}

\emph{R1 for a justification of either kind --- the mean exceeding the median,
or the long tail on the right. Naming the shape with no reason earns the A1
only. ``Skewed left'' earns nothing, however it is justified. FT: award the R1
for a comparison that is correct for the candidate's own mean and median.}

\textbf{(e)} \textbf{Hungary} at $27\%$: it sits well above the rest, with a
gap between it and the next value ($25\%$), far from where the bulk of the data
lies. \hfill \textbf{R1}

\emph{R1 requires the country (or the value 27) \emph{and} a reason of the
right kind --- distance from the rest of the data, or the gap. ``It is the
biggest'' earns nothing: every data set has a biggest value. No $1.5 \times$ IQR
test is expected here; the definition in use is informal.}

\newpage

%--- Problem 2: US state sales tax, left skew [8 marks] ---
\item[2.] [Maximum mark: 8]

Thirteen state sales tax rates, sorted:
$$0,\; 0,\; 0,\; 2.9,\; 4.0,\; 4.0,\; \underline{4.0},\; 4.0,\; 5.3,\; 6.0,\; 6.25,\; 7.0,\; 7.25$$
$n = 13$, $\sum x = 50.7$.

\textbf{(a)} $\bar{x} = \dfrac{50.7}{13} = 3.9$ exactly \hfill \textbf{M1A1}

\emph{M1 for a sum of the thirteen values divided by 13; A1 for 3.9. The
commonest error is to drop the three zeros and divide by 10, giving 5.07: that
earns the M1 only. A state charging $0\%$ is still a data value.}

\textbf{(b)} With $n = 13$ the median is the 7th value: $4.0$
\hfill \textbf{M1A1}

\emph{M1 for an ordered list or for identifying the 7th position; A1 for 4.0.}

\textbf{(c)} Mode $= 4.0$ (Alabama, Georgia, New York and Wyoming).
\hfill \textbf{A1}

\emph{Four states share $4.0$; three share $0$. Do not accept 0, and do not
accept the two quoted together --- the frequencies are not equal.}

\textbf{(d)} The distribution is \textbf{skewed left} (negative skew).
\hfill \textbf{A1}

Because $\bar{x} = 3.9 < \text{median} = 4.0$, and the dot plot has a tail
running down to the three states at $0$ while most values sit between 4 and
7.25. \hfill \textbf{R1}

\emph{The contrast with Problem 1 is the point of the pair: same three
statistics, opposite order, opposite shape. FT as in 1(d).}

\textbf{(e)} Either answer earns the mark if it is argued.
\hfill \textbf{R1}

\emph{Accept ``outliers'', because they sit far below everything else and
distance from the rest of the data is what makes a value an outlier. Accept
``a separate group'', because there are three of them at exactly the same
value and those states have chosen to have no sales tax at all, which is a
policy rather than a stray measurement. Accept ``both''. The mark is for the
reasoning; a bare label earns nothing. This is the open interpretation item.}

\bigskip
\hrule
\bigskip

\textbf{Teacher notes --- not for the student sheet}

\emph{\textbf{Local taxes on top.} The Problem 2 figures are state rates only.
Cities and counties add their own: the population-weighted national average
combined rate is about $7.5\%$, and Louisiana averages about $10.1\%$ combined
against a $5.0\%$ state rate. Alaska appears at $0\%$ here but its boroughs do
levy local sales tax --- a good example of a data value that is true as recorded
and misleading as read.}

\emph{\textbf{How the thirteen states were chosen, and what that costs.} The
selection is pedagogical, not statistical --- a convenience sample picked to
produce a clean left tail, a single mode, an odd $n$ and a sum that divides
exactly. It is not representative, and the numbers say so:}

{\small
\begin{center}
\renewcommand{\arraystretch}{1.2}
\begin{tabular}{|l|c|c|c|c|c|}
\hline
 & $n$ & Mean & Median & Mode & \% at $0$ \\ \hline
The thirteen on the sheet & 13 & 3.90 & 4.00 & 4.0 & 23\% \\ \hline
The other 37 states       & 37 & 5.51 & 6.00 & 6.0 & 5\% \\ \hline
\textbf{All 50 states}    & 50 & \textbf{5.09} & \textbf{6.00} & \textbf{6.0} & 10\% \\ \hline
\end{tabular}
\end{center}}

\emph{The sheet's mean sits $1.2$ points below the true 50-state mean and its
median a full $2$ points below. Three of the five zero-rate states are in
(Alaska and New Hampshire are not), so zeros are over-represented more than
twofold. The real mode is $6.0\%$ --- ten states charge it --- and only one of
those ten (Florida) is on the sheet. Eight of the thirteen rank in the bottom
twenty of fifty; nothing sits between rank 5 and rank 15.}

\emph{\textbf{What survives.} The conclusion does. All fifty states are also
skewed left, and more sharply than the sheet: mean $-$ median is $-0.91$ across
all fifty against $-0.10$ for the thirteen. So the shape students find is the
real shape; only the statistics are shifted. That is the briefing, and it ties
straight back to 1.2 --- this data set is exactly the kind of purposive sample
they were asked to criticise last week, and it still gets the qualitative answer
right. Worth asking them why a biased sample can preserve a shape while moving
every number.}

\emph{\textbf{City income tax, for the discussion about putting two levels on
one graph.} New York City residents pay a city income tax on top of New York
State's, at four rates from $3.078\%$ to $3.876\%$ (single filers: $3.078\%$ on
the first \$12{,}000 of city taxable income, then $3.762\%$, $3.819\%$ and
$3.876\%$ above \$50{,}000). Yonkers levies a surcharge as well. Most US cities
levy none at all. Two ways to put that on a graph are worth raising with the
class: stack the city rate on the state rate as a bar with two segments, so the
total is the bar height and the split is visible; or plot state rates alone and
mark the few cities separately, since a variable that is zero for almost every
city is not really the same variable as the state rate. The first answers ``what
does a resident pay'', the second answers ``how do states differ'' --- which
question you are asking decides the graph.}

\end{enumerate}
\end{document}
